\section{Introduction}
% ===============================================================

\subsection{From counting roots to counting zeros and poles}

The Fundamental Theorem of Algebra says that a nonzero polynomial $P \in \CC[z]$
of degree $d$ has exactly $d$ roots, counted with multiplicity. This is a
counting statement: the local data (a root, together with an integer recording
how many times it is a root) adds up to a global invariant (the degree).

Polynomials, however, are a very restricted class of functions. As soon as we
allow quotients --- rational functions such as $(z-1)^2/z$ --- two things change.
First, the function may blow up: at $z = 0$ our example has a \emph{pole}, which
we should think of as a zero of negative multiplicity. Second, and less
obviously, the count above stops balancing unless we look at $\infty$ as well.
Our example has a double zero at $z=1$ and a simple pole at $z=0$, and the two
multiplicities $2$ and $-1$ do not cancel. The missing contribution sits at
infinity: as $|z| \to \infty$ the function grows like $z$, and once we make sense
of ``the order of $f$ at $\infty$'' it will turn out to be $-1$. The three
numbers $2$, $-1$, $-1$ now sum to zero.

This is the phenomenon the paper is about. To state it we need a space on which
$\infty$ is an honest point --- the complex projective line $\PP$, equivalently
the Riemann sphere --- and a bookkeeping device that records all the local
multiplicities at once. That device is the divisor.

\subsection{The three objects, informally}

Before the formal definitions, here is what the three main objects mean.

\begin{itemize}[leftmargin=1.5em,itemsep=0.25em]
\item A \emph{divisor} is a finite formal sum $D = \sum_p n_p\, p$ of points of
$\PP$ with integer coefficients. It is nothing more than a way of writing down a
finite list of points together with an integer attached to each. Divisors add
coefficientwise, so they form an abelian group $\Div(\PP)$. The divisor
$\divv(f)$ of a function records its zeros with positive coefficients and its
poles with negative ones.
\item A divisor is \emph{principal} if it is $\divv(f)$ for some nonzero rational
function $f$ --- that is, if the prescribed configuration of zeros and poles is
actually achieved by a function. These form a subgroup $\Prin(\PP)$.
\item The \emph{degree} $\deg(D) = \sum_p n_p$ adds up all the coefficients. It
is the total signed multiplicity, and it is a homomorphism $\Div(\PP) \to \ZZ$.
\end{itemize}

The question the paper answers is: \emph{which} divisors are principal? Since not
every list of points and integers can plausibly come from a function, one expects
some obstruction. The content of the main theorem is that on $\PP$ there is
exactly one obstruction, and it is the degree.

\subsection{Statement of results}

\begin{namedtheorem}\label{thmA}
On $\PP$ the principal divisors are exactly the divisors of degree zero:
\[
\Prin(\PP) \;=\; \ker\big(\deg : \Div(\PP) \to \ZZ\big).
\]
\end{namedtheorem}

This is proved as Theorem~\ref{thm:main} in Section~\ref{sec:main}. Two
consequences follow.

\begin{namedtheorem}\label{thmB}
The degree induces an isomorphism $\Cl(\PP) := \Div(\PP)/\Prin(\PP)
\xrightarrow{\ \sim\ } \ZZ$.
\end{namedtheorem}

\begin{namedtheorem}\label{thmC}
The assignment $D \mapsto \Ostr(D)$ induces an isomorphism $\Cl(\PP) \cong
\Pic(\PP)$, and consequently $\Pic(\PP) \cong \ZZ$: every holomorphic line bundle
on $\PP$ is isomorphic to $\Ostr(n)$ for a unique $n \in \ZZ$.
\end{namedtheorem}

Theorem~\ref{thmB} is Theorem~\ref{thm:cl} and Theorem~\ref{thmC} is
Theorem~\ref{thm:pic}. We emphasise that Theorem~\ref{thmC} is proved here
directly, by exhibiting line bundles through their transition functions on the
two standard charts, rather than by quoting the general divisor--line-bundle
correspondence; see Remark~\ref{rem:algebraic} for the relation to the algebraic
formulation.

\subsection{Prerequisites and outline}

We assume the basic theory of holomorphic functions of one variable (power series,
Laurent expansion, the maximum modulus principle, Cauchy estimates, and the
existence of primitives with vanishing periods), as in Ahlfors \cite{Ahlfors},
together with elementary group theory through the First Isomorphism Theorem. No
prior acquaintance with Riemann surfaces or algebraic geometry is assumed.

Section~\ref{sec:prelim} collects the local analytic facts: every zero or pole of
a meromorphic function is measured by a unique integer. Section~\ref{sec:sphere}
constructs $\PP$ and proves that its meromorphic functions are exactly the
rational ones. Section~\ref{sec:divisors} assembles the local integers into
divisors. Section~\ref{sec:exact} records the two exact sequences that organise
the material. Section~\ref{sec:main} proves Theorem~\ref{thmA},
Section~\ref{sec:cl} proves Theorem~\ref{thmB}, and Section~\ref{sec:pic} proves
Theorem~\ref{thmC}.

% ===============================================================
